\section{Isometric enlargements of the physical space}

% The results in this section are project-derived coordinate statements.  The
% cited CPSV16 passages are construction context, not source statements of the
% invariance theorem below.

\begin{definition}[Doubled physical matrix]
    \label{def:mpdo_doubled_physical_matrix}
    \lean{MPOTensor.doubledPhysicalMatrix}
    \leanok
    Let $V:\C^d\to\C^e$ be a linear map.  Its action on the doubled
    ket--bra index is the matrix $W:\C^{d^2}\to\C^{e^2}$ with
    \begin{align}
        W_{(i,j),(p,q)}
         & =V_{i,p}\,\overline{V_{j,q}}.
        \label{eq:mpdo_doubled_physical_matrix}
    \end{align}
\end{definition}

\begin{theorem}[Transport under an isometric physical embedding]
    \label{thm:mpdo_isometric_physical_embedding}
    \lean{MPSTensor.transferMap_kraus_isometry,
        MPSTensor.isLeftCanonical_kraus_isometry,
        MPSTensor.isInjective_kraus_isometry,
        MPSTensor.isInjective_of_kraus_mixing_isInjective,
        MPSTensor.isInjective_kraus_isometry_iff,
        MPSTensor.GaugeEquiv.sum_smul,
        MPOTensor.doubledPhysicalMatrix_isometry,
        MPOTensor.toMPSTensor_changePhysicalBasis,
        MPOTensor.transferMap_toMPSTensor_changePhysicalBasis,
        MPOTensor.isInjective_toMPSTensor_changePhysicalBasis_iff,
        MPOTensor.isLeftCanonical_toMPSTensor_changePhysicalBasis,
        MPOTensor.isNormalTensor_toMPSTensor_changePhysicalBasis_iff,
        MPOTensor.IsMPDO.changePhysicalBasis,
        MPOTensor.isMPDO_changePhysicalBasis_iff,
        MPOTensor.IsSAL.changePhysicalBasis_of_isometry,
        MPOTensor.isSAL_changePhysicalBasis_iff,
        MPOTensor.physTraceTransfer_changePhysicalBasis,
        MPOTensor.gaugeEquiv_toMPSTensor_changePhysicalBasis,
        MPOTensor.physicalSliceColumnCoisometry,
        MPOTensor.physicalSliceColumnCoisometry_mul_conjTranspose,
        MPOTensor.physicalSliceColumns_changePhysicalBasis,
        MPOTensor.physicalSliceColumns_mul_conjTranspose_changePhysicalBasis,
        MPOTensor.physicalSupportProj_changePhysicalBasis}
    \leanok
    \uses{def:mpdo_doubled_physical_matrix, def:injective, def:left_canonical_tensor, def:nt_cpsv, def:gauge_equiv, def:mpdo, def:is_sal}
    Let $\mathcal K=\{\mathcal K^{pq}\}_{p,q=0}^{d-1}$ be an MPO tensor and
    set
    $\mathcal K_V^{ij}=\sum_{p,q}V_{i,p}\overline{V_{j,q}}\mathcal K^{pq}$.
    Write $A^{(p,q)}=\mathcal K^{pq}$ and
    $A_V^{(i,j)}=\mathcal K_V^{ij}$ for the associated doubled-index MPS
    tensors.  For every $V$, the matrix in
    (\ref{eq:mpdo_doubled_physical_matrix}) satisfies
    $A_V^\tau=\sum_\sigma W_{\tau,\sigma}A^\sigma$.

    Suppose now that $V^\dagger V=\Id$. Then
    \begin{align}
        W^\dagger W               & =\Id,
                                  & \E_{A_V}                                                      & =\E_A,
        \label{eq:mpdo_embedding_transfer}                                                                                      \\
        A_V\text{ is injective}
                                  & \Longleftrightarrow A\text{ is injective},
                                  & A_V\text{ is normal}
                                  & \Longleftrightarrow A\text{ is normal},
        \notag                                                                                                                  \\
        \mathcal K_V\text{ is an MPDO}
                                  & \Longleftrightarrow \mathcal K\text{ is an MPDO},
                                  & \mathcal K_V\text{ saturates the area law}
                                  & \Longleftrightarrow \mathcal K\text{ saturates the area law},
        \notag                                                                                                                  \\
        \mathcal T_{\mathcal K_V} & =\mathcal T_{\mathcal K},
                                  & \Pi_{\mathcal K_V}                                            & =V \Pi_{\mathcal K}V^\dagger.
        \label{eq:mpdo_embedding_trace_support}
    \end{align}
    Here $\Pi_{\mathcal K}$ is the orthogonal projection onto the joint column
    span of the physical slices of $\mathcal K$.  Left-canonical normalization
    of $A$ implies left-canonical normalization of $A_V$.

    For arbitrary $V$, positivity of every positive-length periodic operator
    for $\mathcal K$ implies the same positivity for $\mathcal K_V$.
    Moreover, if two doubled-index tensors are related by one virtual gauge,
    applying the same $V$ to their physical indices preserves that gauge.

    This theorem is project-derived and is not stated in
    \cite{Cirac2016MPDO_arXiv}.  Its application is to the normal-block and
    per-sector constructions in lines~217--246, 1628--1665, and 1740--1782 of
    that source.
\end{theorem}

\begin{proof}
    \leanok
    The $(\sigma,\rho)$ entry of $W^\dagger W$ factors into the product of
    the corresponding ket and bra entries of $V^\dagger V$, which gives the
    first identity in (\ref{eq:mpdo_embedding_transfer}).  Isometric mixing of
    a Kraus family therefore gives the second identity.  Conversely, each
    original matrix is recovered from the mixed family by
    $A^\sigma=\sum_\tau\overline{W_{\tau,\sigma}}A_V^\tau$; this proves the
    injectivity equivalence.  The transfer-map identity and the equivalence
    between irreducibility of a Kraus family and of its transfer map preserve
    the three clauses of normality without a new spectral argument.

    At length $N$, the periodic operator of $\mathcal K_V$ is
    $V^{\otimes N}\rho^{(N)}(\mathcal K)(V^{\otimes N})^\dagger$.
    This proves positivity for arbitrary $V$.  Under the isometry hypothesis,
    traces and the mutual informations of every contiguous block agree in the
    two physical spaces, which gives the MPDO and area-law equivalences in
    (\ref{eq:mpdo_embedding_trace_support}).  Contracting the two physical
    indices and using $V^\dagger V=\Id$ gives the equality of the matrices
    $\mathcal T$.

    Finally, let $Z_{\mathcal K}$ stack all columns of all physical slices.
    Directly from the definition,
    $Z_{\mathcal K_V}=VZ_{\mathcal K}R$, where $R$ is the blockwise copy of
    $V^\dagger$ and $RR^\dagger=\Id$.  Hence
    $Z_{\mathcal K_V}Z_{\mathcal K_V}^\dagger
        =VZ_{\mathcal K}Z_{\mathcal K}^\dagger V^\dagger$,
    and taking column-support projections proves the last identity in
    (\ref{eq:mpdo_embedding_trace_support}).  A common virtual similarity
    distributes termwise through the physical sums, proving the gauge
    assertion.
\end{proof}
